% القسم: المناهج
% الفصل: الاستقراء
% المبحث: التعريفات الاستقرائية

\documentclass[../../../include/open-logic-section]{subfiles}

\begin{document}

\olfileid{mth}{ind}{idf}

\olsection{التعريفات الاستقرائية}

كثيراً جداً ما نعرّف في المنطق أنواعاً من الأشياء \emph{استقرائياً}،
أي بوضع قواعد تبين ما الذي يعد شيئاً من النوع المراد تعريفه، وتوضح كيفية
الحصول على أشياء جديدة من ذلك النوع انطلاقاً من أشياء قديمة منه. فكثيراً
ما نعرّف، مثلاً، أنواعاً خاصة من متواليات الرموز، كحدود لغة ما
و!!{formula} فيها، بالاستقراء. ولننظر، كمثال بسيط، في السلاسل المؤلفة من
الحروف $\mathrm{a}$ و$\mathrm{b}$ و$\mathrm{c}$ و$\mathrm{d}$، والرمز
~$\circ$، والقوسين $[$ و$]$، مثل
«$[[\mathrm{c} \circ \mathrm{d}][$» و«$[\mathrm{a}[]\circ]$»
و«$\mathrm{a}$» و«$[[\mathrm{a} \circ \mathrm{b}]\circ
\mathrm{d}]$». ولعلك تشعر أن ثمة شيئاً «خاطئاً» في السلسلتين الأوليين:
فالأقواس غير «متوازنة» إطلاقاً في الأولى، وقد تشعر أن على «$\circ$» أن
«يربط» عبارتين لهما معنى في ذاتيهما. أما السلسلتان الثالثة والرابعة
فتبدوان أفضل: فلكل «$[$» قوس إغلاق «$]$»، إن وجدت أقواس أصلاً، ولكل
$\circ$ نستطيع إيجاد عبارتين «سليمتين» على جانبيه تحيط بهما زوج من
الأقواس.

نريد أن نحدد بدقة ما الذي يعد «حداً سليماً». أولاً، كل حرف بمفرده سليم.
وكل ما لا يقتصر على حرف بمفرده ينبغي أن يكون على الصورة
«$[t \circ s]$»، حيث $s$ و$t$ سليمان في ذاتيهما. وبالعكس، إذا كان
$t$ و$s$ سليمين، استطعنا تكوين حد سليم جديد بوضع $\circ$ بينهما
وإحاطتهما بقوسين. ويمكن أن نستعمل هذه العمليات في \emph{تعريف} مجموعة
الحدود السليمة. وهذا \emph{تعريف استقرائي}.

\begin{defn}[الحدود السليمة]
  تعرّف مجموعة \emph{الحدود السليمة} استقرائياً كما يأتي:
  \begin{enumerate}
  \item كل حرف من $\mathrm{a}$ و$\mathrm{b}$ و$\mathrm{c}$
    و$\mathrm{d}$ حد سليم.
  \item إذا كان $s_1$ و$s_2$ حدين سليمين، فإن
    $[s_1 \circ s_2]$ حد سليم أيضاً.
  \item لا شيء آخر حد سليم.
  \end{enumerate}
\end{defn}

يخبرنا هذا التعريف أن شيئاً ما يعد حداً سليماً إذا وفقط إذا أمكن بناؤه
بحسب الشرطين (1) و(2) في عدد متناهٍ من الخطوات. ففي الخطوة الأولى نبني
جميع الحدود السليمة المؤلفة من حروف منفردة، أي:
\[
\mathrm{a}, \mathrm{b}, \mathrm{c}, \mathrm{d}
\]
وفي الخطوة الثانية نطبق (2) على الحدود التي بنيناها، فنحصل على:
\[
[\mathrm{a} \circ \mathrm{a}], [\mathrm{a} \circ \mathrm{b}],
[\mathrm{b} \circ \mathrm{a}], \dots, [\mathrm{d} \circ \mathrm{d}]
\]
لجميع توليفات حرفين. وفي الخطوة الثالثة نطبق (2) مرة أخرى على أي حدين
سليمين بنيناهما حتى ذلك الحين. فنحصل على حدود سليمة جديدة مثل
$[\mathrm{a} \circ [\mathrm{a} \circ \mathrm{a}]]$---حيث $t$ هو
$\mathrm{a}$ من الخطوة~1 و$s$ هو
$[\mathrm{a} \circ \mathrm{a}]$ من الخطوة~2---ومثل
$[[\mathrm{b} \circ \mathrm{c}] \circ [\mathrm{d} \circ
\mathrm{b}]]$، المبني من الحدين $[\mathrm{b} \circ \mathrm{c}]$
و$[\mathrm{d} \circ \mathrm{b}]$ من الخطوة~2. وهكذا. ويمنع البند (3)
تسلل أي شيء لم يُبن بهذه الطريقة إلى مجموعة الحدود السليمة.

لاحظ أننا لم نثبت بعد أن كل متوالية رموز «تبدو» سليمة هي سليمة بحسب
هذا التعريف. لكن ينبغي أن يكون واضحاً أن كل ما نستطيع بناءه «يبدو
سليماً» بالفعل: فالأقواس متوازنة، و$\circ$ يربط جزأين سليمين في
ذاتيهما.

الميزة الأساسية للتعريفات الاستقرائية هي أن التعريف يخبرك، إذا أردت
إثبات شيء عن جميع الحدود السليمة، بالحالات التي يجب النظر فيها. فإذا
قيل لك، مثلاً، إن $t$ حد سليم، يخبرك التعريف الاستقرائي بالصورة التي
يمكن أن يكون عليها $t$: فإما أن يكون حرفاً، وإما أن يكون
$[s_1 \circ s_2]$ لبعض حدين سليمين $s_1$ و~$s_2$. وبسبب البند (3)،
فهذان هما الاحتمالان الوحيدان.

عند إثبات قضايا عن جميع عناصر مجموعة معرّفة استقرائياً، تكتسب الصورة
القوية للاستقراء أهمية خاصة. افترض، مثلاً، أننا نريد إثبات أن عدد
الرموز $[$ في كل حد سليم طوله~$n$ هو~$< n/2$. ويمكن النظر إلى هذا
بوصفه قضية عن جميع~$n$: لكل $n$، عدد الرموز $[$ في أي حد سليم
طوله~$n$ هو $< n/2$.

\begin{prop}
  لكل $n$، عدد الرموز $[$ في حد سليم طوله~$n$ هو $< n/2$.
\end{prop}

\begin{proof}
لإثبات هذه النتيجة بالاستقراء القوي، علينا أن نبين صدق القضية الشرطية
الآتية:
\begin{quote}
  إذا كان كل حد سليم طوله~$l$ يحتوي على عدد من الرموز $[$ أقل من
  $l/2$ لكل $l < k$، فإن كل حد سليم طوله~$k$ يحتوي على عدد منها أقل
  من $k/2$.
\end{quote}
ولإثبات هذه القضية الشرطية، افترض صدق مقدمها؛ أي افترض أن الحدود السليمة
ذات الطول~$l$ تحتوي على عدد من الرموز $[$ أقل من $l/2$ لأي $l<k$.
ونسمي هذا الافتراض فرض الاستقراء. ونريد أن نبين صدق الشيء نفسه للحدود
السليمة ذات الطول~$k$.

افترض إذن أن $t$ حد سليم طوله~$k$. ولأن الحدود السليمة معرّفة
استقرائياً، فلدينا حالتان: (1)~أن يكون $t$ حرفاً بمفرده، أو (2)~أن يكون
$t$ هو $[s_1 \circ s_2]$ لبعض حدين سليمين $s_1$ و~$s_2$.
\begin{enumerate}
\item $t$ حرف. عندئذ $k = 1$، وعدد الرموز $[$ في $t$ هو~$0$.
وبما أن $0 < 1/2$، تصح القضية.
\item $t$ هو $[s_1 \circ s_2]$ لبعض حدين سليمين $s_1$ و~$s_2$.
  ليكن $l_1$ طول~$s_1$ و$l_2$ طول~$s_2$. عندئذ يكون طول~$t$، وهو
  $k$، مساوياً $l_1+l_2+3$، أي مجموع طولي $s_1$ و$s_2$ مع الرموز
  الثلاثة $[$ و$\circ$ و$]$. ولما كان $l_1+l_2+3$ أكبر دائماً من
  $l_1$، فإن $l_1 < k$. وبالمثل $l_2 < k$. وهذا يعني أن فرض
  الاستقراء ينطبق على الحدين $s_1$ و~$s_2$: فعدد الرموز $[$ في
  $s_1$، وليكن~$m_1$، أقل من $l_1/2$، وعددها في~$s_2$، وليكن~$m_2$،
  أقل من $l_2/2$.

  عدد الرموز $[$ في $t$ هو عددها في~$s_1$ مضافاً إليه عددها في~$s_2$
  ثم~$1$، أي إنه $m_1 + m_2 + 1$. وبما أن $m_1 < l_1/2$ و
  $m_2 < l_2/2$، فلدينا:
  \[
  m_1 + m_2 + 1 < \frac{l_1}{2} + \frac{l_2}{2} + 1 = \frac{l_1+l_2+2}{2} < \frac{l_1+l_2+3}{2} = k/2.
  \]
\end{enumerate}
لقد بينا، في كلتا الحالتين، أن عدد الرموز $[$ في $t$ أقل من $k/2$،
بالاستناد إلى فرض الاستقراء. فتتبع القضية بالاستقراء القوي.
\end{proof}

\begin{prob}
  عرّف مجموعة الحدود فائقة السلامة كما يأتي:
  \begin{enumerate}
  \item كل حرف من $\mathrm{a}$ و$\mathrm{b}$ و$\mathrm{c}$
    و$\mathrm{d}$ حد فائق السلامة.
  \item إذا كان $s$ حداً فائق السلامة، فإن $[s]$ كذلك.
  \item إذا كان $s_1$ و$s_2$ حدين فائقي السلامة، فإن
    $[s_1 \circ s_2]$ كذلك.
  \item لا شيء آخر حد فائق السلامة.
  \end{enumerate}
  بيّن أن عدد الرموز $[$ في حد فائق السلامة~$t$ طوله~$n$ لا يزيد على
  $n/2 +1$.
\end{prob}

\end{document}
